% ===================================================================
%  圖表第 1 號 — 學校。--- 中學。--- 課室。
%
%  Ceci n'est PAS une transcription : c'est une traduction, faite en
%  2026, en CANTONAIS ECRIT d'aujourd'hui, en caracteres
%  traditionnels. D'ou trois differences avec texto/io et texto/fr,
%  et trois seulement :
%
%   * pas de \nl ni de \cc — la traduction ne suit aucune ligne
%     imprimee, et les lignes de ce fichier ne sont que du confort ;
%   * pas de \begin{VUpage} — elle n'a ni feuillet ni folio, et la
%     page de lecture ne lui en affiche pas ;
%   * le gras se pose sur le TERME seul, ponctuation dehors.
%
%  Tout le reste est identique, et doit l'etre : les cles « %%K » sont
%  celles de texto/io, macro pour macro pour l'apparat, et les renvois
%  \textsuperscript{(n)} visent les MEMES objets de la planche.
%
%  LE CANTONAIS ECRIT N'EST PAS DU CHINOIS STANDARD LU A LA
%  CANTONAISE, ET TOUTE LA COLONNE EN DEPEND. C'est meme la seule
%  raison pour laquelle elle existe : le WU, qui a autant de
%  locuteurs, a ete ecarte parce que qui le parle ecrit le mandarin,
%  de sorte qu'une colonne « wuu » aurait double « zh » a la lettre
%  pres. Le cantonais, lui, a ses caracteres a lui, et Hong Kong les
%  imprime tous les jours. Si cette colonne ecrivait 是, 在, 的, 不,
%  了, 他, elle serait la colonne chinoise avec une autre police —
%  et elle n'aurait pas lieu d'etre.
%
%  ELLE ECRIT DONC, PARTOUT, LES FORMES CANTONAISES :
%
%    係 (et non 是)      être              喺 (et non 在)      à, dans
%    嘅 (et non 的)      de, attributif    唔 (et non 不)      ne pas
%    咗 (et non 了)      accompli          佢 (et non 他)      il, elle
%    哋 (pluriel)        我哋, 佢哋        啲                  les, du
%    呢 / 嗰             ce-ci / ce-là     睇 (et non 看)      regarder
%    冇 (et non 沒有)    ne pas avoir      畀 (et non 給)      donner
%    攞 (et non 拿)      prendre           好 (et non 很)      très
%
%  Ces quatorze mots reviennent des centaines de fois dans les seize
%  tableaux : ce sont eux, et non le vocabulaire, qui font la colonne.
%
%  LE VOCABULAIRE SCOLAIRE, LUI, EST CELUI DE HONG KONG, et il differe
%  de celui de Pekin plus souvent qu'on ne croit :
%
%    課室    la classe        (Pekin : 教室)
%    老師    le maitre        (Pekin : 教員, 老师)
%    功課    le devoir        (Pekin : 作业)
%    擦膠    la gomme         (Pekin : 橡皮)
%    書包    le cartable      — le meme, mais avec 個 pour classificateur
%    黑板擦  la brosse        (Pekin : 板擦)
%
%  LES NOMS DE PERSONNE SONT CEUX DE LA MISSION DE HONG KONG, et ils
%  ne sont pas ceux de Pekin : ils ont ete transcrits sur la
%  prononciation CANTONAISE des missionnaires, un siecle avant que le
%  mandarin ne serve de norme. 約翰 se lit ici « joek3 hon6 » et non
%  « yuehan ». Les latins de la note (*) ne bougent pas : c'est d'eux
%  qu'elle parle.
%
%  « postigo » (t01, renvoi 78) EST UN VASISTAS, ET LA PLANCHE LE DIT.
%  Verifie sur la gravure en 2026, non sur les dictionnaires : la
%  verriere de la classe a deux rangees de carreaux, et dans la rangee
%  HAUTE deux petits vantaux sont graves battants, gonds sur un
%  montant vertical, ouverts vers l'interieur. C'est le seul ouvrant
%  de la verriere, et il ne sert qu'a aerer. Le cantonais dit
%  « 氣窗 », la fenetre a air, qui nomme exactement cela.
% ===================================================================

%%K t01-apar-0 apar
\VUsaut{2.0mm}
\VUcentre{12.6pt}{120}{說明手冊}
\VUsaut{3.2mm}
\VUcentre{8.4pt}{160}{附於}
\VUsaut{2.6mm}
\VUtitre{15.6pt}{0}{德爾馬斯輔助掛圖}
\VUsaut{3.4mm}
\VUornamento{9pt}
\VUsaut{5.0mm}
\VUcentre{13.2pt}{90}{第一組}
\VUsaut{3.0mm}
\VUfilet{16mm}
\VUsaut{5.6mm}

%%K t01-apar-1 apar
\VUtitre{15.0pt}{60}{圖表第 1 號}
\VUsaut{1.4mm}
\VUtitre{11.4pt}{0}{學校。--- 中學。--- 課室。}
\VUsaut{2.2mm}
\VUfilet{20mm}
\VUsaut{6.4mm}

%%K t01-tit-1 sub
\VUpk{10.2pt}{40}{總體描述。}
\VUsaut{3.0mm}

%%K t01-01-1 p
1. --- 呢張圖表畫嘅係一間\VUgras{中學}嘅\VUgras{課室}。呢間課室入面有八張
\VUgras{檯} \textsuperscript{(18)} 同八張\VUgras{長凳} \textsuperscript{(32)}。每張
長凳坐三個學生。\VUgras{老師} \textsuperscript{(7)} 坐喺\VUgras{椅}
\textsuperscript{(8)} 上面，喺自己嘅\VUgras{講檯} \textsuperscript{(6)} 度。

%%K t01-02-1 p
2. --- 講檯左邊有個\VUgras{玻璃門櫃} \textsuperscript{(15)}。右邊係\VUgras{黑板}
\textsuperscript{(1)}，仲有\VUgras{黑板擦} \textsuperscript{(2)} 同\VUgras{粉筆}
\textsuperscript{(3)}。

%%K t01-02-2 p
講檯上面掛住幾幅\VUgras{掛圖} \textsuperscript{(9, 11, 12)} 同一幅\VUgras{地圖}
\textsuperscript{(10)}。

%%K t01-03-1 p
3. --- 唔係個個學生都坐喺自己嘅\VUgras{位} \textsuperscript{(17)} 度。約翰
\textsuperscript{(4)} 企喺黑板旁邊；佢攞住黑板擦，抹緊啲\VUgras{幾何圖形}。

%%K t01-03-2 p
彼得喺講檯附近；佢收緊啲\VUgras{功課} \textsuperscript{(5)}，攞去畀老師。

%%K t01-04-1 p
4. --- \VUgras{呢位}老師係教\VUgras{現代語}嘅。佢教啲學生講、讀同寫外國話
\VUgras{（德文、英文、西班牙文、意大利文、俄文}或者\VUgras{法文）}。

%%K t01-05-1 p
5. --- 課室好大好光。入面盡頭有一扇闊嘅\VUgras{窗}，上面有兩個\VUgras{氣窗}
\textsuperscript{(78)}，仲有一道\VUgras{玻璃門}，用嚟通風同採光。

%%K t01-05-2 p
呢間課室唔凍。中間有個好好嘅\VUgras{火爐} \textsuperscript{(46)} 同佢嘅
\VUgras{煙通} \textsuperscript{(45)} 用嚟取暖。

%%K t01-05-3 p
牆上釘住幾個\VUgras{衣鈎} \textsuperscript{(58)}；學生嘅帽同大褸就掛喺嗰度。

%%K t01-tit-2 sub
\VUsaut{5.2mm}
\VUpk{10.2pt}{40}{操場。}
\VUsaut{3.6mm}

%%K t01-06-1 p
6. --- \VUgras{鐘} \textsuperscript{(63)} 指住九點零五分。

%%K t01-06-2 p
\VUgras{小息} \textsuperscript{(76)} 啱啱完，堂又開始返。小息係喺\VUgras{操場}
\textsuperscript{(75)} 度玩嘅。我哋見到有幾個學生喺度玩。一位\VUgras{訓導}
\textsuperscript{(74)} 睇住佢哋。

%%K t01-07-1 p
7. --- 喺操場度我哋仲見到第啲人：\VUgras{督學} \textsuperscript{(73)}、
\VUgras{校長} \textsuperscript{(71)} 同\VUgras{教務主任} \textsuperscript{(72)}。

%%K t01-07-2 p
督學視察啲學校，校長管住間中學，教務主任就睇住啲功課。

%%K t01-07-3 p
第四個人係\VUgras{看更} \textsuperscript{(77)}。佢挾住本簿，用嚟記低邊個缺席。

%%K t01-08-1 p
8. --- 操場盡頭有\VUgras{體操器械} \textsuperscript{(70)} 同做\VUgras{手工}
\textsuperscript{(69)} 嘅工場。

%%K t01-08-2 p
圖表右邊我哋開始見到\VUgras{音樂}同\VUgras{美術}嘅\VUgras{課室}。

%%K t01-noto-1 noto
\VUnotes{143.2mm}{%
(*) 講到\textit{教名}，我哋勸大家連拉丁文嘅寫法都一齊抄低。咁做係避免無數
異寫嘅唯一辦法。況且，\textit{Giovanni、Jean、}\textit{Johann、Jan、Hans、John、
Ivan}呢啲，點揀好呢？難道要為教名編一本專門嘅字典？我哋不如喺拉丁文攞佢哋嘅
主格，就咁講：\VUgras{Ioannes、Ioanna；Iulius、Iulia；Iakobus；Andreas；
Lukas。}（完整文法）。%
}

%%K t01-tit-3 sub
\VUpk{10.2pt}{40}{詳細描述。}
\VUsaut{2.4mm}

%%K t01-tit-4 sub
\VUpk{10.2pt}{40}{好學生。}
\VUsaut{3.0mm}

%%K t01-09-1 p
9. --- 講檯右邊第一張長凳嘅學生係\VUgras{亨利}(*)、\VUgras{司提反}同
\VUgras{查理}。

%%K t01-09-2 p
亨利好專心好勤力；佢聽住老師講。佢張檯上面有一本\VUgras{簿}
\textsuperscript{(28)}、一本\VUgras{書} \textsuperscript{(29)}、一本\VUgras{字典}
\textsuperscript{(30)} 同一個\VUgras{筆盒} \textsuperscript{(31)}。佢本書有好多
\VUgras{頁}；每一章有幾個\VUgras{段落}，每一\VUgras{行}有好多\VUgras{字}。

%%K t01-09-4 p
佢隔籬嘅司提反 \textsuperscript{(24)} 好用功。佢喺簿上面抄低下堂要用嘅嘢。佢寫得
一手好\VUgras{字}。佢本書閂咗，我哋淨係見到個\VUgras{書皮} \textsuperscript{(27)}。
佢個\VUgras{圓規} \textsuperscript{(26)} 就喺佢隔籬。

%%K t01-09-5 p
查理 \textsuperscript{(23)} 攞住本書；佢揭開嚟溫書。同每個學生一樣，佢面前有個
\VUgras{墨水樽} \textsuperscript{(25)}。佢好聽話。

%%K t01-10-1 p
10. --- 隔籬張檯坐住\VUgras{路易}、\VUgras{安東}同\VUgras{保羅}。路易背緊佢嘅
\VUgras{書} \textsuperscript{(22)}，又答老師嘅\VUgras{問題}。安東
\textsuperscript{(21)} 好唔專心，老師有時仲要罰佢添。保羅 \textsuperscript{(20)}
係個好同學，又和氣又肯幫人。佢好專心。

%%K t01-tit-5 sub
\VUsaut{4.2mm}
\VUpk{10.2pt}{40}{壞學生。}
\VUsaut{3.2mm}

%%K t01-11-1 p
11. --- 亨利後面坐住\VUgras{佐治} \textsuperscript{(38)}。佢唔算係曳仔，不過佢
\VUgras{多口}、唔專心、坐唔定。佢嘅\VUgras{輕浮}同\VUgras{唔聽話}喺堂上搞出
\VUgras{混亂}，所以佢成日要俾人\VUgras{嚴厲警告}。佢本\VUgras{草稿簿}
\textsuperscript{(35)} 好污糟，佢用支\VUgras{筆嘴} \textsuperscript{(34)} 喺上面
\VUgras{搞到成篇墨漬}，所以佢成日要用\VUgras{擦膠} \textsuperscript{(36)}；佢個
\VUgras{書袋} \textsuperscript{(33)} 爛咗，因為佢好大意。點解佢會攞住個
\VUgras{筆筒} \textsuperscript{(37)} 入現代語堂呢？

%%K t01-11-2 p
佢隔籬嘅愛德蒙 \textsuperscript{(39)} 唔中意佢。呢個雖然的確有啲懶，但係佢靜靜哋
唔搞事。佢唔多口；不過佢唔聽書，寧願睇自己本\VUgras{讀本}入面嘅\VUgras{故事}。

%%K t01-11-3 p
呢張長凳第三個學生係\VUgras{路多維}\textsuperscript{(40)}。佢好勤力。佢喺一張白
\VUgras{紙}上面寫嘢。佢面前擺住張\VUgras{吸墨紙} \textsuperscript{(41)}。

%%K t01-12-1 p
12. --- 老師左邊第二張長凳嘅學生仲細個。第一個係細細個嘅\VUgras{愛美}，佢仲未
識寫得好；佢帶住自己嘅\VUgras{石板} \textsuperscript{(42)}、\VUgras{石筆}同
\VUgras{海綿} \textsuperscript{(43)}。

%%K t01-12-2 p
佢隔籬係\VUgras{奧古}同\VUgras{伯納} \textsuperscript{(44)}。後者讀書幾好，不過
佢嘅\VUgras{衣著}好邋遢。

%%K t01-13-1 p
13. --- 佢哋後面坐住三個\VUgras{懶蟲。亞爾拔}唔溫書。\VUgras{雅各}
\textsuperscript{(47)} 唔聽書，喺度瞓覺。佢嘅\VUgras{懶惰}為佢招嚟\VUgras{責備}
甚至\VUgras{懲罰}。最後，\VUgras{基士頓} \textsuperscript{(48)} 太過輕浮。佢
\VUgras{品行}唔好，又完全唔肯改。

%%K t01-14-1 p
14. --- 相反，坐喺佢隔籬嘅同學品行好好。\VUgras{恩尼斯} \textsuperscript{(51)}
中意\VUgras{整齊}同\VUgras{有規矩}；佢成日用自己支\VUgras{間尺}
\textsuperscript{(50)} 同支\VUgras{鉛筆} \textsuperscript{(49)}。佢係\VUgras{走讀生}。

%%K t01-14-2 p
\VUgras{方濟} \textsuperscript{(52)} 係\VUgras{寄宿生}；佢著校服，喺學校食飯瞓覺。
佢係個好\VUgras{同學}。

%%K t01-14-3 p
毛里斯用自己把\VUgras{小刀} \textsuperscript{(56)} 削\VUgras{鉛筆}；佢好細心。

%%K t01-14-4 p
佢帶齊自己所有書，佢本\VUgras{文法} \textsuperscript{(55)}、佢本\VUgras{記事簿}
\textsuperscript{(54)}，仲有塊\VUgras{墊寫板} \textsuperscript{(53)}。佢個
\VUgras{書包} \textsuperscript{(57)} 就喺佢後面地下。

%%K t01-14-5 p
課室盡頭嗰兩張檯係\VUgras{好學生} \textsuperscript{(19)} 嘅位。

%%K t01-tit-6 sub
\VUsaut{4.6mm}
\VUpk{10.2pt}{40}{美術同音樂。}
\VUsaut{3.4mm}

%%K t01-15-1 p
15. --- \VUgras{美術室}入面有幾個靚\VUgras{模型}掛喺牆上，仲有一盞有\VUgras{燈罩}
嘅\VUgras{燈} \textsuperscript{(62)} 吊喺天花。\VUgras{老師} \textsuperscript{(60)}
好有耐性；佢改學生嘅\VUgras{畫} \textsuperscript{(61)}，又教佢哋照真人畫
\VUgras{半身像} \textsuperscript{(59)}。

%%K t01-16-1 p
16. --- \VUgras{音樂室}好似好有趣。喺嗰度人哋學讀\VUgras{樂譜}，唱寫喺
\VUgras{五線譜}上面嘅\VUgras{音階音} \textsuperscript{(67)}（do、re、mi、fa、sol、
la、si\,=\,C. D. E. F. G. A. B.），學識啲\VUgras{譜號}，又唱好聽嘅\VUgras{曲調}。
啲樂譜擺喺\VUgras{譜架} \textsuperscript{(65)} 上面。一個學生\VUgras{彈鋼琴}
\textsuperscript{(64)}，\VUgras{音樂老師} \textsuperscript{(66)} 就\VUgras{打拍子}
\textsuperscript{(68)}。

%%K t01-17-1 p
17. --- 我哋開始見到做\VUgras{手工} \textsuperscript{(69)} 嘅\VUgras{工場}
一角，喺嗰度啲男仔可以學刨木同銼鐵。唔係間間中學都有呢啲。

%%K t01-tit-7 sub
\VUsaut{5.0mm}
\VUpk{10.2pt}{40}{科學室。}
\VUsaut{3.6mm}

%%K t01-18-1 p
18. --- 數學老師教學生\VUgras{幾何}、\VUgras{算術}、\VUgras{計算}同\VUgras{米制}。
我哋喺黑板上面見到幾何嘅主要圖形：\VUgras{圓} \textsuperscript{(a)}、
\VUgras{正方形} \textsuperscript{(b)}、\VUgras{三角形} \textsuperscript{(c)}、
\VUgras{線} \textsuperscript{(d)}。

%%K t01-18-2 p
我哋仲見到一條\VUgras{運算}；佢唔係\VUgras{加法}，唔係\VUgras{減法}，亦都唔係
\VUgras{乘法}：佢正正係\VUgras{除法} \textsuperscript{(e)}。

%%K t01-19-1 p
19. --- 喺米制嗰幅掛圖上面，我哋見到：\VUgras{米} \textsuperscript{(a)}、
\VUgras{天秤} \textsuperscript{(b)} 同佢啲\VUgras{砝碼}、\VUgras{公升}
\textsuperscript{(e)}、\VUgras{十升} \textsuperscript{(f)}、\VUgras{分升}同
\VUgras{立方米} \textsuperscript{(d)}。

%%K t01-19-2 p
學生仲學\VUgras{自然科學} \textsuperscript{(11)} --- 動物學 \textsuperscript{(a)}、
植物學 \textsuperscript{(b)}、地質學 \textsuperscript{(c)} --- 同\VUgras{歷史}
\textsuperscript{(12)}。

%%K t01-19-3 p
櫃入面擺住\VUgras{物理} \textsuperscript{(15)} 同\VUgras{化學} \textsuperscript{(16)}
嘅儀器。

%%K t01-19-4 p
櫃頂擺住：一個\VUgras{球體} \textsuperscript{(14)}、一個\VUgras{圓錐}同一個
\VUgras{地球儀} \textsuperscript{(13)}。

%%K t01-19-5 p
啲掛圖上面掛住一幅\VUgras{地圖}，畫住世界五大洲：\VUgras{歐洲}
\textsuperscript{(g)}、\VUgras{亞洲} \textsuperscript{(e)}、\VUgras{非洲}
\textsuperscript{(f)}、\VUgras{美洲} \textsuperscript{(ab)} 同\VUgras{大洋洲}
\textsuperscript{(h)}，仲有兩個大洋，\VUgras{太平洋} \textsuperscript{(c)} 同
\VUgras{大西洋} \textsuperscript{(d)}。
\VUsaut{6.0mm}
\VUfilet{22mm}
